\parttitle[popecology]{Population and Community Ecology}{Populations, their growth and regulation, and every way species interact.}

\section{The scope of ecology and its levels of organization}

Ecology is the study of the distribution and abundance of organisms and of the
interactions that determine them. The two nouns matter: \tm{distribution} is a
question about space (where is the species, and where is it not), and
\tm{abundance} is a question about number (how many, and why not more). Every
question on the USABO ecology section reduces to one of those two, or to the
mechanism linking them.

\begin{defbox}[title={DEFINITIONS: the levels of ecological organization}]
\nr{44-popecology-levels}{Levels of ecological organization}
\begin{itemize}
\item \tm{Organismal ecology} treats the single individual: its physiology,
behaviour and morphology as adaptations to the abiotic environment. Subdivided
into physiological, evolutionary and behavioural ecology.
\item \tm{Population ecology} treats a group of conspecific individuals living in
one area at one time, with the potential to interbreed. The state variable is
$N$, the number of individuals.
\item \tm{Community ecology} treats all populations of all species in an area and
their interactions. The state variables are species richness $S$ and the
interaction coefficients.
\item \tm{Ecosystem ecology} adds the abiotic compartment and follows energy
(J~m$^{-2}$~yr$^{-1}$) and matter (g~m$^{-2}$~yr$^{-1}$) through it.
\item \tm{Landscape ecology} treats a mosaic of ecosystems, patch geometry and
flows between patches, typically over $1$ to $10^4$~km$^2$.
\item \tm{Global ecology} (the biosphere) treats the whole planet: the carbon
cycle, ocean circulation, and climate feedbacks.
\end{itemize}
\end{defbox}

\nr{44-popecology-population-def}{Population, biological definition}
A population is not merely a set of individuals of one species. The operational
definition requires (i) one species, (ii) one defined area, (iii) one time, and
(iv) actual or potential gene exchange. A wolf pack in Yellowstone and a wolf
pack in Alaska are not one population, because condition (iv) fails at any useful
timescale even though (i) holds.

\begin{thmbox}[title={PRINCIPLE: emergent properties at each level}]
\nr{44-popecology-emergent}{Emergent properties of levels}
Each level has properties undefined at the level below. An individual has an age,
but only a population has an \tm{age structure}. An individual is born or dies,
but only a population has a \tm{birth rate}. A population has a density, but only
a community has a \tm{diversity}. The consequence for the exam: if a question
asks for a rate per capita, it is a population question; if it asks for
J~m$^{-2}$~yr$^{-1}$, it is an ecosystem question.
\end{thmbox}

\begin{tabular}{@{}lp{3.4cm}p{4.6cm}@{}}
\toprule
\textbf{Level} & \textbf{State variable} & \textbf{Characteristic question}\\
\midrule
Organism & Body mass, rate & How does it tolerate $-40\degc$?\\
Population & $N$, $r$, $K$ & Why does $N$ stop at $500$?\\
Community & $S$, $H'$, $\alpha_{ij}$ & Why do $12$ warblers coexist?\\
Ecosystem & Energy, nutrient flux & Where do $2000$~kcal~m$^{-2}$ go?\\
Landscape & Patch number, edge & Does a $50$~m corridor restore flow?\\
Biosphere & Global pools & Why is \ce{CO2} at $425$~ppm?\\
\bottomrule
\end{tabular}

\begin{obsbox}[title={WORTH NOTICING: the scales span twenty orders of magnitude}]
\nr{44-popecology-scales}{Spatial and temporal scales in ecology}
A bacterial population responds over $10^{-6}$~m and $20$~min. A forest stand
responds over $10^{2}$~m and $10^{2}$~yr. The biosphere responds over $10^{7}$~m
and $10^{4}$ to $10^{6}$~yr. Because the scales differ by roughly $10^{13}$ in
space and $10^{10}$ in time, a process that is a constant at one scale is a
variable at another: climate is a fixed parameter for a population study of one
season and the dependent variable for a biosphere study.
\end{obsbox}

\begin{warnbox}[title={TRAP: ecology is not environmental science and not environmentalism}]
\nr{44-popecology-vs-envsci}{Ecology against environmental science}
The wrong belief is that ecology means protecting nature. Ecology is a
descriptive and predictive natural science; it makes no value claims.
\tm{Environmental science} is interdisciplinary and applied: it imports ecology,
chemistry, geology, economics, law and policy to manage human impact, and it does
take a normative position (pollution is bad). \tm{Environmentalism} is a social
and political movement. On the exam, a stem asking "which practice should be
adopted" is environmental science; a stem asking "what will $N$ do" is ecology.
A further distinction: ecology includes systems with no human involvement at all,
which environmental science by definition does not.
\end{warnbox}

\begin{tipbox}[title={TECHNIQUE: identify the level before doing any work}]
Trigger: any ecology stem with numbers in it. Read the units first.
\begin{itemize}
\item Units of individuals per area $\rightarrow$ population, use $N$ models.
\item Units of individuals per area per time $\rightarrow$ population rate.
\item Dimensionless index between $0$ and $1$ (or $0$ and $\ln S$)
$\rightarrow$ community diversity.
\item Units of mass or energy per area per time $\rightarrow$ ecosystem
productivity.
\item Ratio of two areas or a species count against area $\rightarrow$ island
biogeography or landscape.
\end{itemize}
Choosing the level chooses the equation, and almost all lost marks in this
section are marks lost by applying an ecosystem formula to a population stem.
\end{tipbox}

\begin{keybox}[title={MEMORIZE: the abiotic and biotic factor split}]
\nr{44-popecology-abiotic-biotic}{Abiotic against biotic factors}
\tm{Abiotic} factors are non-living: temperature, water, light
($\approx 1360$~W~m$^{-2}$ at the top of the atmosphere, $\approx 1000$~W~m$^{-2}$
at noon at sea level), salinity ($35$~ppt in open ocean, $0.5$~ppt fresh),
pH, soil texture, fire, wind, dissolved \ce{O2} ($\approx 8$--$10$~mg~L$^{-1}$
in cool fresh water, under $2$~mg~L$^{-1}$ is hypoxic). \tm{Biotic} factors are
living: competitors, predators, prey, parasites, mutualists, pathogens. The rule
that matters: abiotic factors act mostly in a \tm{density independent} way, biotic
factors mostly in a \tm{density dependent} way, and that difference drives the
whole regulation debate.
\end{keybox}

\section{Population properties: density, dispersion, sex ratio and age structure}

A population is described by a small set of measurable properties. Four are
examinable in detail, and each has a distinct method of measurement.

\begin{keybox}[title={MEMORIZE: density and its two flavours}]
\nr{44-popecology-density}{Population density}
\[
D=\frac{N}{A}
\]
$N$ = number of individuals, $A$ = area (or volume, for plankton and soil fauna).
Units: individuals~m$^{-2}$, ha$^{-1}$ or km$^{-2}$.
\nr{44-popecology-crude-ecological}{Crude against ecological density}
\tm{Crude density} uses total area. \tm{Ecological density} (specific density)
uses only the habitable area. A lake of $100$~ha holding $4000$ trout gives a
crude density of $40$~ha$^{-1}$; if only the $25$~ha littoral zone is usable, the
ecological density is $160$~ha$^{-1}$, a fourfold difference. Ecological density
is the biologically meaningful number and the one that correlates with
competition.
\end{keybox}

\nr{44-popecology-density-ranges}{Typical density ranges}
Densities span an enormous range and the exam likes order-of-magnitude checks:
soil bacteria $10^{8}$--$10^{10}$~g$^{-1}$; soil nematodes $10^{6}$~m$^{-2}$;
earthworms $100$--$500$~m$^{-2}$ in good pasture; small rodents
$10$--$200$~ha$^{-1}$; white-tailed deer $4$--$20$~km$^{-2}$; African elephant
$0.5$--$2$~km$^{-2}$; tiger $0.5$--$5$ per $100$~km$^{-2}$. Density scales roughly
as $M^{-0.75}$ with body mass $M$, the mirror of the metabolic
$\frac{3}{4}$~power law.

\begin{defbox}[title={DEFINITIONS: the three dispersion patterns}]
\nr{44-popecology-dispersion}{Dispersion patterns}
\tm{Dispersion} (also called spacing or pattern) is the arrangement of
individuals in space. Do not confuse it with \tm{dispersal}, which is movement.
\begin{itemize}
\item \tm{Clumped} (aggregated, contagious): individuals in patches. Caused by
patchy resources, social behaviour, limited seed dispersal, or clonal growth.
This is by far the commonest pattern in nature, quoted at $70$--$90\%$ of
surveyed species.
\item \tm{Uniform} (regular, even): individuals more evenly spaced than chance.
Caused by antagonism: territoriality in birds, allelopathy in creosote bush,
self-thinning in plantations, nesting gannets at $\approx 2$ nests~m$^{-2}$.
\item \tm{Random}: position of one individual independent of all others. Requires
a homogeneous environment and no interaction. Rare; the classic example is
wind-dispersed forest spores and some clams on uniform mudflats.
\end{itemize}
\end{defbox}

The pattern is diagnosed statistically, not by eye, and the statistic is the
variance-to-mean ratio of counts in equal-sized quadrats. For a random pattern
the counts follow a Poisson distribution, in which the variance equals the mean.

\begin{calcbox}[title={WORKED: diagnosing dispersion from the variance-to-mean ratio}]
\nr{44-popecology-variance-mean}{Variance-to-mean ratio (index of dispersion)}
Ten $1$~m$^2$ quadrats give seedling counts $0,0,1,0,7,0,9,1,0,2$.
\ttitle{Mean.} $\bar{x}=(0+0+1+0+7+0+9+1+0+2)/10=20/10=2.0$.
\ttitle{Sum of squared deviations.} Deviations are $-2,-2,-1,-2,5,-2,7,-1,-2,0$;
squares $4,4,1,4,25,4,49,1,4,0$, summing to $96$.
\ttitle{Variance.} $s^2=96/(10-1)=10.67$.
\ttitle{Index.} $s^2/\bar{x}=10.67/2.0=5.33$.
\[
\frac{s^2}{\bar{x}}=5.33\gg 1 \quad\Rightarrow\quad \text{clumped}
\]
\ttitle{Significance test.} $\chi^2=(n-1)s^2/\bar{x}=9\times 5.33=48.0$ on
$9$~df; the $0.05$ critical value is $16.9$, so the departure from randomness is
significant.
\ttitle{Check.} Six of ten quadrats are empty while two hold $7$ and $9$; an
empty-quadrat fraction of $0.60$ against the Poisson prediction
$e^{-2}=0.135$ confirms aggregation. Decision rule:
$s^2/\bar{x}\approx 1$ random, $>1$ clumped, $<1$ uniform. \tm{$\surd$}
\end{calcbox}

\begin{warnbox}[title={TRAP: clumped is the default, and scale changes the answer}]
\nr{44-popecology-dispersion-trap}{Scale dependence of dispersion}
Two errors. First, students answer "random" because it sounds like the null; in
nature clumped is the commonest pattern and random is the rarest. Second, and more
subtle, dispersion is scale dependent: the same forest sampled with $1$~m$^2$
quadrats may look clumped (trees in groves), with $100$~m$^2$ quadrats may look
random, and with $1$~cm$^2$ quadrats looks uniform because two stems cannot
occupy the same point. If a question changes quadrat size and asks why the
conclusion changed, scale dependence is the answer.
\end{warnbox}

\begin{defbox}[title={DEFINITIONS: sex ratio and age structure}]
\nr{44-popecology-sex-ratio}{Sex ratio}
\tm{Sex ratio} is conventionally males per $100$ females, or the proportion male.
\tm{Primary} sex ratio is at conception, \tm{secondary} at birth, \tm{tertiary}
at maturity. In humans the secondary ratio is $\approx 105$ males per $100$
females ($1.05$), the tertiary ratio falls below $1.00$ by age $50$ and to
$\approx 0.45$ at age $85$ because male mortality is higher at every age.
Fisher's principle predicts a $1{:}1$ investment ratio at equilibrium; deviations
signal unequal cost or local mate competition.
\nr{44-popecology-operational-sex-ratio}{Operational sex ratio}
The \tm{operational sex ratio} counts only sexually receptive adults and is the
number that predicts the intensity of sexual selection; in elephant seals it can
exceed $1{:}40$ in effective terms because one bull monopolizes $40$ or more cows.
\nr{44-popecology-age-structure}{Age structure}
\tm{Age structure} is the proportion of the population in each age class, split
into \tm{pre-reproductive}, \tm{reproductive} and \tm{post-reproductive} classes.
It is displayed as an \tm{age pyramid} with age on the vertical axis and the two
sexes back to back.
\end{defbox}

\begin{tabular}{@{}lp{3.0cm}p{4.4cm}@{}}
\toprule
\textbf{Pyramid shape} & \textbf{Growth} & \textbf{Signature and example}\\
\midrule
Broad triangle & Rapid growth & Wide pre-reproductive base, $>40\%$ under $15$; Niger, Afghanistan\\
Bell or dome & Slow or stable & Even classes, $15$--$20\%$ under $15$; USA, France\\
Urn or inverted & Declining & Narrow base, $>25\%$ over $65$; Japan, Italy, Germany\\
\bottomrule
\end{tabular}

\begin{obsbox}[title={WORTH NOTICING: momentum hidden in the pyramid}]
\nr{44-popecology-momentum}{Population momentum}
A population whose total fertility rate falls to the replacement value of
$2.1$ today keeps growing for $40$ to $70$ years, because the broad
pre-reproductive base has yet to pass through the reproductive classes. This is
\tm{population momentum}, and it is why China's population continued to rise for
decades after fertility fell below $2.1$ around $1992$. The exam version: a
stable age structure is a consequence of constant vital rates, not a cause, and
it takes roughly one generation time to establish.
\end{obsbox}

\begin{tipbox}[title={TECHNIQUE: matching the property to the method}]
For each property, know the one method that measures it.
\begin{itemize}
\item Density of sessile organisms: quadrat counts, direct.
\item Density of mobile animals: mark-recapture, or transect with a detection
correction.
\item Density of secretive species: indirect indices (pellet counts, tracks,
calls, camera-trap rates, eDNA), which give relative not absolute density.
\item Dispersion: variance-to-mean ratio, or nearest-neighbour distance.
\item Sex ratio: sexed captures or harvest records, corrected for
sex-biased catchability.
\item Age structure: growth rings in scales, otoliths, teeth or horns; tooth wear
in ungulates; skull suture closure; or mark-recapture of known-age cohorts.
\end{itemize}
\end{tipbox}

\section{Sampling methods: quadrats, transects, mark-recapture and the Lincoln index}

You almost never count every individual. A \tm{total count} (census) is possible
only for large conspicuous organisms in small areas: nesting colonies from aerial
photographs, elephants from the air, trees over $10$~cm diameter in a plot. Every
other estimate is a sample, and the exam tests whether you know what each method
assumes.

\begin{defbox}[title={DEFINITIONS: the plot and line methods}]
\nr{44-popecology-quadrat}{Quadrat sampling}
A \tm{quadrat} is a bounded sample plot, usually square, placed at random or on a
grid. Count or estimate cover inside, then scale up:
$\hat{N}=\bar{x}\times(A_{\text{total}}/a_{\text{quadrat}})$. Quadrat size is
chosen so that most quadrats contain some individuals but few contain very many:
$0.25$~m$^2$ for grassland herbs, $1$~m$^2$ for tall herbs, $4$--$25$~m$^2$ for
shrubs, $100$--$1000$~m$^2$ for trees.
\nr{44-popecology-transect}{Transects}
A \tm{line transect} records every individual touching a line; a \tm{belt
transect} is a strip of fixed width (commonly $0.5$--$2$~m) treated as a series
of contiguous quadrats; a \tm{point transect} takes records at fixed intervals
along the line. Transects are the method of choice when the variable of interest
changes along a gradient (shore height, altitude, distance from a pollution
source) because they capture \tm{zonation} that random quadrats would average
away.
\nr{44-popecology-cover-frequency}{Cover, frequency, abundance}
\tm{Percentage cover} is the fraction of ground the vertical projection of a
species covers, and may total over $100\%$ where canopies overlap.
\tm{Frequency} is the percentage of quadrats in which a species occurs, and is
strongly quadrat-size dependent. \tm{Density} is counts per area and needs
countable individuals, so it fails for clonal grasses.
\end{defbox}

\begin{tabular}{@{}lp{3.2cm}p{4.4cm}@{}}
\toprule
\textbf{Method} & \textbf{Best for} & \textbf{Main limitation}\\
\midrule
Random quadrats & Sessile, evenly spread & Misses gradients; many needed if clumped\\
Belt transect & Zonation on a gradient & Not random, so no simple variance\\
Line intercept & Cover of shrubs, corals & Poor for small sparse species\\
Mark-recapture & Mobile animals & Needs $4$ strong assumptions\\
Removal trapping & Small mammals, fish & Destructive; assumes constant effort\\
Indices (pellets, calls) & Rare or shy species & Relative only, needs calibration\\
\bottomrule
\end{tabular}

\begin{keybox}[title={MEMORIZE: the Lincoln-Petersen index and its relatives}]
\nr{44-popecology-lincoln}{Lincoln-Petersen index}
Mark $M$ individuals, release, allow mixing, then take a second sample of $n$
individuals of which $m$ are marked. Assuming the marked fraction in the sample
equals the marked fraction in the population,
\[
\frac{M}{\hat{N}}=\frac{m}{n}\qquad\Longrightarrow\qquad
\hat{N}=\frac{Mn}{m}
\]
\nr{44-popecology-chapman}{Chapman estimator}
The Lincoln index is biased upward for small samples and undefined when $m=0$.
The bias-corrected \tm{Chapman estimator} is preferred whenever $m<10$:
\[
\hat{N}_{C}=\frac{(M+1)(n+1)}{m+1}-1
\]
\nr{44-popecology-lincoln-se}{Standard error of the Lincoln estimate}
\[
\mathrm{SE}(\hat{N})\approx\sqrt{\frac{M^2 n (n-m)}{m^3}}
\]
and the $95\%$ interval is $\hat{N}\pm 1.96\,\mathrm{SE}$.
\nr{44-popecology-schnabel}{Schnabel method}
For $k>2$ sampling occasions,
$\hat{N}=\sum(C_t M_t)/\sum R_t$ where $C_t$ is the catch, $M_t$ the marked
animals at large and $R_t$ the recaptures on occasion $t$. It is far more precise
because it pools information.
\end{keybox}

\begin{calcbox}[title={WORKED: mark-recapture of a beetle population}]
A researcher traps $120$ ground beetles, marks each with a dot of enamel on the
elytra, and releases them. Four days later a second trapping session yields $150$
beetles, of which $18$ carry marks.
\ttitle{Lincoln estimate.}
\[
\hat{N}=\frac{Mn}{m}=\frac{120\times 150}{18}=\frac{18000}{18}=1000
\]
\ttitle{Chapman estimate.}
\[
\hat{N}_C=\frac{(121)(151)}{19}-1=\frac{18271}{19}-1=961.6-1=960.6\approx 961
\]
\ttitle{Standard error.}
\[
\mathrm{SE}=\sqrt{\frac{120^2\times 150\times(150-18)}{18^3}}
=\sqrt{\frac{14400\times150\times132}{5832}}
\]
$14400\times150=2\,160\,000$; $\times 132=285\,120\,000$; divided by $5832$ gives
$48\,888$; the square root is $221$.
\ttitle{Interval.} $1000\pm 1.96\times 221 = 1000\pm 433$, so roughly
$570$ to $1430$. The estimate is imprecise because only $18$ recaptures were
obtained.
\ttitle{Density.} If the trapping grid covered $0.50$~ha, the density is
$1000/0.50=2000$~beetles~ha$^{-1}$, i.e. $0.2$~m$^{-2}$.
\ttitle{Check.} The marked fraction of the second sample is
$18/150=0.12$, and $120/1000=0.12$ as well, so the proportions match by
construction. The Chapman value is slightly lower than the Lincoln value, which
is the expected direction of the bias correction. \tm{$\surd$}
\end{calcbox}

\begin{warnbox}[title={TRAP: the four assumptions of mark-recapture}]
\nr{44-popecology-mr-assumptions}{Assumptions of mark-recapture}
The wrong belief is that $\hat{N}=Mn/m$ is a measurement. It is a model estimate
resting on four assumptions, each with a predictable direction of error.
\begin{itemize}
\item \tm{Closed population}: no birth, death, immigration or emigration between
samples. Any dilution of marks by births or immigrants inflates $\hat{N}$.
\item \tm{Equal catchability}: marked and unmarked animals equally likely to be
caught. If marking makes animals \tm{trap-shy}, $m$ falls and $\hat{N}$ is
overestimated; if it makes them \tm{trap-happy} (bait learning), $m$ rises and
$\hat{N}$ is underestimated.
\item \tm{Marks retained and readable}: lost marks reduce $m$ and inflate
$\hat{N}$.
\item \tm{Complete mixing}: enough time for marked animals to redistribute, but
not so long that assumption one fails. This is the tension that sets the
interval between samples.
\end{itemize}
A fifth practical rule: the mark must not raise mortality or predation risk, so
paint on a moth's wing or a collar on a bird must be small relative to body mass
(under $3$--$5\%$ is the usual ethical ceiling).
\end{warnbox}

\begin{expbox}[title={EXPERIMENT: Petersen's plaice, 1896, and Lincoln's ducks, 1930}]
\nr{44-popecology-petersen-lincoln}{Origin of the mark-recapture estimator}
C.~G.~J.~Petersen tagged plaice in the Limfjord with numbered bone buttons and
used the recovery ratio in the commercial catch to estimate stock size and
migration; this is the first quantitative use of the ratio estimator.
F.~C.~Lincoln in $1930$ applied the same arithmetic at continental scale to North
American waterfowl, using the ratio of returned leg bands to hunter-shot birds to
estimate the total duck population, which is why the estimator carries his name in
North American texts and Petersen's in European ones. What the work established:
a single ratio, applied to a marked subsample, converts a catch into a population
size, provided the marked and unmarked animals are mixed and equally vulnerable.
Both authors noted the equal-catchability problem, which is still the main
weakness a century later.
\end{expbox}

\begin{tipbox}[title={TECHNIQUE: reading a mark-recapture stem for the trick}]
Trigger: numbers $M$, $n$, $m$ appear. Before computing, scan for the sabotage.
\begin{itemize}
\item "Two weeks later" in a species that breeds weekly $\rightarrow$ closure
violated, estimate too high.
\item "Marked with a bright tag" and a visual predator present $\rightarrow$
differential mortality, estimate too high.
\item "Trapped at the same baited stations" $\rightarrow$ trap-happy, estimate too
low.
\item "Marks faded" $\rightarrow$ $m$ too low, estimate too high.
\item $m$ small (under $10$) $\rightarrow$ use Chapman, and expect a wide
interval.
\end{itemize}
\end{tipbox}

\section{Life tables, survivorship curves and fecundity schedules}

A life table is an age-specific accounting of mortality and reproduction. It
turns a population from a single number $N$ into a schedule, and from the
schedule you can compute every growth parameter.

\begin{keybox}[title={MEMORIZE: the life table columns}]
\nr{44-popecology-lifetable}{Life table notation}
\begin{itemize}
\item $x$ = age or age class (years, instars, stages).
\item $n_x$ = number alive at the start of age $x$.
\item $l_x=n_x/n_0$ = \tm{survivorship}, the proportion of the original cohort
alive at age $x$. Always $l_0=1$.
\item $d_x=n_x-n_{x+1}$ = number dying during interval $x$.
\item $q_x=d_x/n_x$ = \tm{age-specific mortality rate}, the probability of dying
during $x$ given survival to $x$. This is the column that reveals the risky ages.
\item $p_x=1-q_x$ = age-specific survival probability.
\item $e_x$ = \tm{life expectancy} remaining at age $x$, computed as
$e_x=\sum_{y\ge x}L_y/n_x$ with $L_y=(n_y+n_{y+1})/2$.
\item $m_x$ (or $b_x$) = \tm{fecundity}, mean number of female offspring per
surviving female of age $x$. The fecundity schedule is the list of $m_x$.
\end{itemize}
\end{keybox}

\begin{defbox}[title={DEFINITIONS: cohort, static and segment life tables}]
\nr{44-popecology-cohort-static}{Cohort against static life tables}
A \tm{cohort} (horizontal, dynamic) life table follows a single group born at the
same time until all are dead. It gives the true $l_x$ but requires a study longer
than the lifespan, so it is feasible for annual plants, insects and barnacles and
impractical for oaks or elephants.
A \tm{static} (vertical, time-specific) life table takes a snapshot of the age
distribution of the living population, or a sample of ages at death, and infers
$l_x$ from it. It is quick but assumes the population has a \tm{stable age
distribution} and constant vital rates. If those fail, the table is an artefact
of past good and bad years.
A \tm{segment} life table follows several cohorts over a few age classes each and
splices the results.
\end{defbox}

\begin{calcbox}[title={WORKED: completing a life table and computing $R_0$}]
\nr{44-popecology-R0}{Net reproductive rate $R_0$}
A cohort of $n_0=500$ annual plants gives the following.

\begin{tabular}{@{}rrrrrr@{}}
\toprule
$x$ & $n_x$ & $l_x$ & $d_x$ & $q_x$ & $m_x$\\
\midrule
0 & 500 & 1.000 & 400 & 0.800 & 0\\
1 & 100 & 0.200 & 60 & 0.600 & 2.0\\
2 & 40 & 0.080 & 30 & 0.750 & 5.0\\
3 & 10 & 0.020 & 10 & 1.000 & 8.0\\
4 & 0 & 0.000 & --- & --- & ---\\
\bottomrule
\end{tabular}

\ttitle{Column arithmetic.} $l_1=100/500=0.200$; $d_0=500-100=400$;
$q_0=400/500=0.800$. Note $q_2=30/40=0.750$ exceeds $q_1=0.600$, so mortality
rises again after age $1$ even though fewer individuals die in absolute terms.
\ttitle{Net reproductive rate.}
\[
R_0=\sum_x l_x m_x = (1.000)(0)+(0.200)(2.0)+(0.080)(5.0)+(0.020)(8.0)
\]
$=0+0.400+0.400+0.160=0.960$.
\ttitle{Generation time.}
\[
T=\frac{\sum x\, l_x m_x}{\sum l_x m_x}
=\frac{(1)(0.400)+(2)(0.400)+(3)(0.160)}{0.960}
=\frac{0.400+0.800+0.480}{0.960}=\frac{1.680}{0.960}=1.75\ \text{yr}
\]
\ttitle{Intrinsic rate of increase, approximated.}
\[
r\approx\frac{\ln R_0}{T}=\frac{\ln 0.960}{1.75}=\frac{-0.0408}{1.75}
=-0.0233\ \mathrm{yr^{-1}}
\]
\ttitle{Check.} $R_0=0.96<1$ means each female is replaced by $0.96$ daughters,
so the population must shrink, and indeed $r<0$. The two results agree in sign,
which is the check to run every time: $R_0>1\Leftrightarrow r>0$;
$R_0=1\Leftrightarrow r=0$; $R_0<1\Leftrightarrow r<0$. The decline is about
$2.3\%$ per year, so the population halves in
$\ln 2/0.0233\approx 30$~yr. \tm{$\surd$}
\end{calcbox}

\begin{thmbox}[title={PRINCIPLE: the three survivorship curves}]
\nr{44-popecology-survivorship}{Survivorship curve types}
Plot $\log_{10} l_x$ (or $\ln l_x$) against age, conventionally starting from
$1000$ survivors.
\begin{itemize}
\item \tm{Type I} (convex): low juvenile mortality, mortality concentrated in old
age. Large mammals, humans, elephants. Few, large, well-provisioned offspring;
human $l_x$ stays above $0.9$ to age $50$ in a developed country.
\item \tm{Type II} (straight on a log axis): constant $q_x$, so a constant
\emph{probability} of death at all ages. Hydra, many birds after fledging,
squirrels, some lizards, and annual plants in the adult phase.
\item \tm{Type III} (concave): massive early mortality then low adult mortality.
Oysters, most marine invertebrates, oaks, most fish. An oyster releasing
$10^{8}$ eggs may see $l_x$ fall below $10^{-6}$ within weeks.
\end{itemize}
A straight line on a log survivorship plot is the diagnostic of Type II because
$\log l_x=\log l_0 - kx$ corresponds to $l_x=e^{-kx}$, a constant hazard.
\end{thmbox}

\begin{warnbox}[title={TRAP: survivorship curves need a logarithmic axis}]
\nr{44-popecology-survivorship-trap}{Log axis on survivorship curves}
The wrong belief is that Type II is a straight line on any axis. Plot $l_x$ on a
linear axis and an exponential decay is a curve, so a genuine Type II population
looks Type III. Only on a log axis does constant proportional mortality become a
straight line. Two further traps: (i) the survivorship curve is not the age
pyramid, since one is a cohort's fate through time and the other is a snapshot of
the living; (ii) most real species are intermediate, and many switch type with
age, for example a plant with Type III seedlings and Type I adults.
\end{warnbox}

\begin{expbox}[title={EXPERIMENT: Murie's Dall sheep skulls, 1944, and Deevey's reanalysis}]
\nr{44-popecology-dall-sheep}{Dall mountain sheep life table}
Adolph Murie collected $608$ skulls of \sn{Ovis dalli} in Denali and aged each
from horn annuli; E.~S.~Deevey converted the age-at-death distribution into a
static life table in $1947$. The result: heavy mortality in the first year
($q_0\approx 0.20$), a long low-mortality plateau from ages $2$ to $8$ with
$q_x$ near $0.03$--$0.06$, and a sharp rise after age $9$, with almost none
surviving past $14$. What it established: (i) a sample of skulls, which is far
easier to obtain than a followed cohort, can yield a life table, provided the
population is stationary and skulls of all ages are equally likely to be found;
(ii) wild large mammals show a broadly Type I curve, refuting the older view that
wild animals die at a constant rate; (iii) wolf predation fell disproportionately
on lambs and on sheep past age $9$, the same classes with high $q_x$, which is
evidence that predation removed animals already near death.
\end{expbox}

\begin{obsbox}[title={WORTH NOTICING: fecundity and survivorship trade off}]
\nr{44-popecology-fecundity-tradeoff}{Fecundity-survivorship trade-off}
Across species the product $l_x m_x$ summed over life is pinned near $1$ in any
persistent population, which forces a trade-off: a species with
$l_{\text{juvenile}}\approx 10^{-6}$ must produce $\approx 10^{6}$ eggs, and a
species producing one offspring at a time must keep juvenile survival above
$\approx 0.5$. An ocean sunfish laying $3\times 10^{8}$ eggs and an albatross
laying one egg every two years both average $R_0\approx 1$ over the long run.
The exam version: if a stem gives you a huge clutch, expect Type III, and if it
gives you parental care, expect Type I.
\end{obsbox}

\section{Exponential growth: the model, $r$, and doubling time}

When resources are unlimited and the environment is constant, each individual
contributes the same per capita surplus, and the population grows in proportion
to its own size. That is exponential growth, the ecological analogue of compound
interest.

\begin{keybox}[title={MEMORIZE: the exponential model in both forms}]
\nr{44-popecology-exponential}{Exponential growth model}
\ttitle{Continuous (overlapping generations).}
\[
\frac{dN}{dt}=(b-d)N=rN,\qquad N_t=N_0e^{rt}
\]
$N$ = population size, $t$ = time, $b$ = instantaneous per capita birth rate,
$d$ = instantaneous per capita death rate, $r=b-d$ = \tm{intrinsic rate of
increase} with units time$^{-1}$ (per individual per unit time, so
individuals cancel).
\ttitle{Discrete (non-overlapping generations).}
\[
N_{t+1}=\lambda N_t,\qquad N_t=N_0\lambda^{t}
\]
$\lambda$ = \tm{finite rate of increase} (geometric growth factor),
dimensionless.
\nr{44-popecology-r-lambda}{Relation between $r$ and $\lambda$}
\[
\lambda=e^{r},\qquad r=\ln\lambda
\]
Growth if $r>0$ or $\lambda>1$; stationary if $r=0$ or $\lambda=1$; decline if
$r<0$ or $0<\lambda<1$. $\lambda$ can never be negative; $r$ can.
\end{keybox}

\begin{defbox}[title={DEFINITIONS: the several rates called $r$}]
\nr{44-popecology-rmax}{Intrinsic rate of increase and biotic potential}
\tm{$r_{\max}$} (also $r_m$, the intrinsic rate of natural increase, or
\tm{biotic potential}) is the value of $r$ realized under ideal conditions:
unlimited resources, optimal temperature, no crowding, no enemies. The
\tm{observed} $r$ is lower by the amount of \tm{environmental resistance}, the
sum of all limiting factors. A useful identity for a closed population:
$r=b-d$; for an open one,
$\Delta N=(B+I)-(D+E)$ with $B$ births, $D$ deaths, $I$ immigrants, $E$
emigrants, and the per capita version is
$r=(b+i)-(d+e)$.
\nr{44-popecology-r-scaling}{Scaling of $r_{\max}$ with body size}
$r_{\max}$ scales as roughly $M^{-0.26}$: \sn{Escherichia coli}
$\approx 60$~day$^{-1}$ under optimum ($20$~min doubling), \sn{Paramecium}
$\approx 1.0$~day$^{-1}$, a flour beetle $\approx 0.1$~day$^{-1}$, a rat
$\approx 0.015$~day$^{-1}$, a human $\approx 0.0003$~day$^{-1}$
($\approx 0.1$~yr$^{-1}$ at maximum), an elephant
$\approx 0.0002$~day$^{-1}$. Small organisms recover from a crash in days, large
ones in decades, and that single fact drives most conservation policy.
\end{defbox}

\begin{calcbox}[title={WORKED: doubling time and projection}]
\nr{44-popecology-doubling}{Doubling time}
Setting $N_t=2N_0$ in $N_t=N_0e^{rt}$ gives $2=e^{rt_2}$, so
\[
t_2=\frac{\ln 2}{r}=\frac{0.693}{r}
\]
\ttitle{Case 1, a bacterium.} \sn{E.~coli} in rich broth at $37\degc$ doubles in
$20$~min, so
$r=0.693/20=0.0347$~min$^{-1}$, i.e. $2.08$~h$^{-1}$ or $49.9$~day$^{-1}$.
Starting from a single cell, after $10$~h ($600$~min):
\[
N=1\times e^{0.0347\times600}=e^{20.8}=1.1\times10^{9}
\]
Consistency check by doublings: $600/20=30$ doublings, $2^{30}=1.07\times10^{9}$.
The two agree to within rounding. \tm{$\surd$}
\ttitle{Case 2, a human population.} A country of $N_0=8.0\times10^{6}$ has a
crude birth rate of $28$ per $1000$~yr$^{-1}$ and a crude death rate of $9$ per
$1000$~yr$^{-1}$, with negligible migration.
\[
r=\frac{28-9}{1000}=0.019\ \mathrm{yr^{-1}}\quad(1.9\%\ \text{per year})
\]
Doubling time $t_2=0.693/0.019=36.5$~yr. The "rule of $70$" shortcut gives
$70/1.9=36.8$~yr, which agrees because $70\approx 100\ln 2$.
Population after $25$~yr:
\[
N=8.0\times10^{6}e^{0.019\times25}=8.0\times10^{6}e^{0.475}
=8.0\times10^{6}\times1.608=1.29\times10^{7}
\]
\ttitle{Check.} $25$~yr is $0.685$ of a doubling time, and
$2^{0.685}=1.609$, matching the factor $1.608$ obtained from the exponential.
The absolute annual increment at the start is
$rN=0.019\times8.0\times10^{6}=152\,000$ per year, and at the end
$0.019\times1.29\times10^{7}=245\,000$ per year: $r$ is constant while the
increment grows, which is the whole point of the model. \tm{$\surd$}
\end{calcbox}

\begin{warnbox}[title={TRAP: $r$, $\lambda$ and percentage growth are three different numbers}]
\nr{44-popecology-r-trap}{Confusing $r$ with $\lambda$ and with percent growth}
The wrong belief is that a population growing $10\%$ per year has $r=0.10$ and
$\lambda=1.10$ interchangeably. If the growth is measured as a discrete annual
census ratio, then $\lambda=1.10$ and $r=\ln 1.10=0.0953$, not $0.10$. The gap
is small at $10\%$ (about $5\%$ relative) but reaches $\lambda=2.0$ against
$r=0.693$ at high rates, a $44\%$ discrepancy. Two further traps: (i) a constant
$r$ does not mean a constant number added, since $dN/dt=rN$ grows with $N$;
(ii) exponential growth plotted on a log axis is a straight line whose slope
\emph{is} $r$, so "straight line on a semi-log plot" means exponential, not
linear.
\end{warnbox}

\begin{expbox}[title={EXPERIMENT: the Protection Island pheasants, 1937--1942}]
\nr{44-popecology-pheasant}{Protection Island pheasant introduction}
Eight ring-necked pheasants ($2$ males, $6$ females) were released on
Protection Island, Washington, a $1.6$~km$^2$ predator-free island with no
resident pheasants. Autumn counts over the next five breeding seasons gave
approximately $40$, $100$, $400$, $770$ and $1740$ birds; the study ended when
the army occupied the island. Fitting $N_t=N_0\lambda^t$ from $8$ to $1740$ over
$5$ years gives
$\lambda=(1740/8)^{1/5}=(217.5)^{0.2}=2.94$, so $r=\ln 2.94=1.08$~yr$^{-1}$ and a
doubling time of $0.64$~yr. What it established: a vertebrate released into an
empty, enemy-free habitat grows very close to geometrically for many
generations, so the exponential model is not a mathematical fiction; and the
realized $\lambda\approx 2.9$ was far above the mainland value near $1.1$,
quantifying how much of the shortfall in ordinary populations is environmental
resistance rather than low biotic potential.
\end{expbox}

\begin{obsbox}[title={WORTH NOTICING: why nothing grows exponentially for long}]
\nr{44-popecology-exp-absurdity}{Absurdity check on exponential growth}
A single \sn{E.~coli} doubling every $20$~min for $48$~h would give
$2^{144}\approx 2\times10^{43}$ cells, about $10^{31}$~kg, roughly $10^{4}$ Earth
masses. A pair of elephants at $r=0.0002$~day$^{-1}$ would, after $750$ years,
cover the continents. Because exponential growth always collides with a finite
resource, an observed exponential phase is a statement about the time window, not
about the species. Real populations show it during colonization, after a
disturbance, in the lag-free phase of a bacterial culture, and in humans from
about $1750$ to $1970$, when $r$ peaked near $0.021$~yr$^{-1}$ in $1965$--$1970$
and has since fallen to about $0.009$~yr$^{-1}$.
\end{obsbox}

\begin{tabular}{@{}lrrp{3.0cm}@{}}
\toprule
\textbf{Organism} & $\boldsymbol{r}$ \textbf{(day$^{-1}$)} & $\boldsymbol{t_2}$ & \textbf{Condition}\\
\midrule
\sn{E.~coli} & 50 & 20 min & Rich broth, $37\degc$\\
\sn{Paramecium} & 1.0 & 17 h & $25\degc$, bacteria fed\\
\sn{Daphnia} & 0.30 & 2.3 d & $20\degc$, algae\\
Flour beetle & 0.10 & 6.9 d & Flour culture\\
Norway rat & 0.015 & 46 d & Ideal, laboratory\\
Human & 0.0003 & 6.3 yr & Maximum recorded\\
\bottomrule
\end{tabular}

\section{Logistic growth: carrying capacity and the sigmoid curve}

The logistic model is the simplest repair to the exponential: make the per capita
growth rate fall linearly as the population fills the environment.

\begin{keybox}[title={MEMORIZE: the logistic equation in three equivalent forms}]
\nr{44-popecology-logistic}{Logistic growth model}
\[
\frac{dN}{dt}=rN\left(\frac{K-N}{K}\right)=rN\left(1-\frac{N}{K}\right)
\]
$K$ = \tm{carrying capacity}, the equilibrium population the environment
sustains, in the same units as $N$. The bracket $(1-N/K)$ is the
\tm{unused portion of the carrying capacity}, running from $1$ at $N\approx 0$ to
$0$ at $N=K$.
\ttitle{Per capita form.}
\[
\frac{1}{N}\frac{dN}{dt}=r\left(1-\frac{N}{K}\right)=r-\frac{r}{K}N
\]
a straight line of intercept $r$ and slope $-r/K$: plotting per capita growth
against $N$ and getting a descending straight line is the signature of logistic
dynamics, and the $x$-intercept is $K$.
\ttitle{Integrated form.}
\[
N_t=\frac{K}{1+\left(\dfrac{K-N_0}{N_0}\right)e^{-rt}}
\]
\end{keybox}

\begin{defbox}[title={DEFINITIONS: carrying capacity and environmental resistance}]
\nr{44-popecology-K}{Carrying capacity}
\tm{Carrying capacity} $K$ is the population size at which births balance
deaths, set by the supply of the single most limiting resource: food, water,
nest sites, territory space, refuges from predation, or the accumulation of
waste. It is a property of the environment and the species jointly, not of the
species alone.
\nr{44-popecology-env-resistance}{Environmental resistance}
\tm{Environmental resistance} is the difference between the exponential
trajectory and the realized one; in the logistic it is exactly the factor
$N/K$ of the potential growth that is lost.
\nr{44-popecology-phases}{Phases of the sigmoid curve}
The \tm{sigmoid} (S-shaped) curve has four phases: a \tm{lag} phase where $N$ is
small and absolute increase is slow, an \tm{exponential} (log) phase where
$N\ll K$ and the bracket is near $1$, a \tm{deceleration} phase past the
inflection where crowding bites, and a \tm{plateau} phase where $N$ oscillates
about $K$.
\end{defbox}

\begin{thmbox}[title={PRINCIPLE: maximum absolute growth occurs at $N=K/2$}]
\nr{44-popecology-inflection}{Inflection point of the logistic at $K/2$}
$dN/dt=rN-rN^2/K$ is a downward parabola in $N$ with roots at $N=0$ and $N=K$,
so its maximum lies midway, at $N=K/2$. Substituting,
\[
\left(\frac{dN}{dt}\right)_{\max}=r\frac{K}{2}\left(1-\frac{1}{2}\right)
=\frac{rK}{4}
\]
This value is the \tm{maximum sustainable yield} (MSY) of the logistic model, and
$N=K/2$ is the stock size a fishery or forestry operation aims to hold. At
$N=K/2$ the curve's inflection occurs, $d^2N/dt^2=0$, and the curve changes from
concave up to concave down.
\nr{44-popecology-msy}{Maximum sustainable yield}
MSY $=rK/4$ at $N^*=K/2$ for the logistic; note that harvesting at exactly MSY
leaves no margin, so any error or environmental downturn drives collapse, which
is why real quotas are set at $0.5$ to $0.8$ of MSY and target stocks above
$K/2$.
\end{thmbox}

\begin{calcbox}[title={WORKED: logistic growth rates and yield for a deer herd}]
A deer herd has $r=0.40$~yr$^{-1}$ and $K=2000$.
\ttitle{Growth rate at four sizes.}
\[
\frac{dN}{dt}=0.40\,N\left(1-\frac{N}{2000}\right)
\]
\begin{tabular}{@{}rrrr@{}}
\toprule
$N$ & $1-N/K$ & $dN/dt$ & Per capita $(1/N)dN/dt$\\
\midrule
100 & 0.950 & 38.0 & 0.380\\
500 & 0.750 & 150.0 & 0.300\\
1000 & 0.500 & 200.0 & 0.200\\
1500 & 0.250 & 150.0 & 0.100\\
1900 & 0.050 & 38.0 & 0.020\\
\bottomrule
\end{tabular}

Arithmetic for $N=1000$: $0.40\times1000=400$; $400\times0.500=200.0$~yr$^{-1}$.
For $N=1500$: $0.40\times1500=600$; $600\times0.250=150.0$~yr$^{-1}$.
\ttitle{Maximum.} $rK/4=0.40\times2000/4=200$~yr$^{-1}$ at $N=1000=K/2$, matching
the table.
\ttitle{Time to reach $N=1000$ from $N_0=100$.} From the integrated form with
$N_t=1000$:
\[
1000=\frac{2000}{1+19e^{-0.40t}}\ \Rightarrow\ 1+19e^{-0.40t}=2.0
\ \Rightarrow\ e^{-0.40t}=\frac{1}{19}
\]
$t=\ln 19/0.40=2.944/0.40=7.36$~yr.
\ttitle{Check, two ways.} First, the table is symmetric: $N=100$ and $N=1900$
both give $38.0$, and $N=500$ and $N=1500$ both give $150.0$, as a parabola
symmetric about $K/2$ must. Second, per capita growth falls linearly from
$0.380$ at $N=100$ to $0.020$ at $N=1900$: the drop is $0.360$ over
$\Delta N=1800$, a slope of $-2.0\times10^{-4}$, which equals
$-r/K=-0.40/2000$ exactly. \tm{$\surd$}
\end{calcbox}

\begin{warnbox}[title={TRAP: what is maximal at $K/2$, and what $K$ is not}]
\nr{44-popecology-logistic-trap}{Logistic misconceptions}
Three errors the exam exploits.
\begin{itemize}
\item \tm{Per capita} growth is maximal at $N\rightarrow 0$, not at $K/2$. Only
the \tm{absolute} rate $dN/dt$ peaks at $K/2$. A stem saying "individuals
reproduce fastest" means per capita and the answer is the smallest $N$.
\item $K$ is not a ceiling the population cannot exceed. Populations overshoot
$K$ routinely, and then $dN/dt$ is negative because $(1-N/K)<0$, so the model
itself predicts decline back toward $K$. A \tm{dieback} after an overshoot is
logistic behaviour with a time lag, not a violation.
\item $K$ is not a constant of the species. It changes with rainfall, season,
fire, and management; a reindeer $K$ on St.~Matthew Island collapsed from
roughly $5000$ to near zero when the lichen mat was eaten out, because the herd
had destroyed the resource that set $K$.
\end{itemize}
\end{warnbox}

\begin{expbox}[title={EXPERIMENT: Gause's \sn{Paramecium} cultures, 1934}]
\nr{44-popecology-gause-logistic}{Gause's logistic cultures}
G.~F.~Gause grew \sn{Paramecium caudatum} in $5$~mL of buffered salt medium with a
fixed daily ration of \sn{Bacillus} cells, counting the whole culture daily.
Numbers rose from $2$ individuals to a plateau near $60$--$65$ per $0.5$~mL
(about $375$ per tube) in $\approx 12$ days, and the daily counts fitted the
logistic curve with $r\approx 1.0$~day$^{-1}$; doubling the bacterial ration
roughly doubled the plateau, while $r$ was unchanged. What it established: $K$
is set by resource supply and $r$ by the organism, that the two parameters are
separable and independently manipulable, and that the sigmoid is not an artefact
of censusing but a real consequence of intraspecific competition for a renewable
ration. Pearl and Reed had already fitted the same curve to the United States
census from $1790$ to $1910$ in $1920$, predicting a plateau near
$197$~million; the actual $2020$ figure of $331$~million shows that a curve
fitting past data does not validate its extrapolation.
\end{expbox}

\begin{tabular}{@{}p{3.4cm}p{3.2cm}p{4.0cm}@{}}
\toprule
\textbf{Assumption} & \textbf{What it requires} & \textbf{How reality breaks it}\\
\midrule
$r$ and $K$ constant & Fixed environment & Seasons, drought change $K$ yearly\\
No time lag & Instant response to density & Gestation, development give lags and overshoot\\
Linear density effect & Each added individual costs the same & Allee effects at low $N$; thresholds at high $N$\\
All individuals equal & No age or size structure & Juveniles neither breed nor compete equally\\
Closed population & No migration & Immigration props up sink populations\\
No genetic change & Fixed vital rates & Selection can raise $r$ or $K$ in tens of generations\\
\bottomrule
\end{tabular}

\begin{obsbox}[title={WORTH NOTICING: the Allee effect inverts the logistic at low $N$}]
\nr{44-popecology-allee}{Allee effect}
The logistic says per capita growth is highest when $N$ is smallest. For social
and sexually reproducing species the opposite holds below a threshold: mates
become hard to find, group defence fails, cooperative hunting breaks down, and
per capita growth falls with falling $N$. This is the \tm{Allee effect}, and a
\tm{strong} Allee effect creates a second equilibrium, an extinction threshold
$A$, below which the population is doomed:
\[
\frac{dN}{dt}=rN\left(1-\frac{N}{K}\right)\left(\frac{N}{A}-1\right)
\]
The passenger pigeon, which nested in colonies of millions and failed to breed in
small flocks, is the standard example; the last wild bird was shot in $1900$ and
the species was extinct by $1914$ despite thousands still alive in $1890$.
Conservation consequence: a minimum viable population is usually quoted in the
range $50$ to $500$ breeding adults for short-term and long-term persistence
respectively, the "$50/500$ rule", now often revised upward to $100/1000$.
\end{obsbox}
